%%%%%%%%%%%%%%%%%%%%%%%%%%%%%%%%%%%%%%%%%%%%%%%%%%%%%%%%%%%%%%%%%%%%%%
% Lớp 8 | Chương III. Tứ giác | Hình bình hành
% Dạng 8H3?12-1: Tính chất hình bình hành
%%%%%%%%%%%%%%%%%%%%%%%%%%%%%%%%%%%%%%%%%%%%%%%%%%%%%%%%%%%%%%%%%%%%%%


% nguồn: De_AI_Toan8_cuoi_hoc_ki_i_1 (nhập 21/11/2026)
\begin{bt}%[8H3H12-1]%[ma:c619aca5]
	Cho hình bình hành $ABCD$ có $\widehat{A} = 110^\circ$, $AB = 8\text{ cm}$ và $BC = 5\text{ cm}$.
a) Tính số đo các góc còn lại của hình bình hành $ABCD$.
b) Tính chu vi của hình bình hành $ABCD$.
	\loigiai{
		a) Vì $ABCD$ là hình bình hành nên các góc đối bằng nhau và hai góc kề một cạnh bù nhau: \\
$\widehat{C} = \widehat{A} = 110^\circ$. \\
$\widehat{B} = 180^\circ - \widehat{A} = 180^\circ - 110^\circ = 70^\circ$. \\
$\widehat{D} = \widehat{B} = 70^\circ$. \\
b) Vì $ABCD$ là hình bình hành nên $CD = AB = 8\text{ cm}$ và $AD = BC = 5\text{ cm}$. \\
Chu vi hình bình hành $ABCD$ là $2(8 + 5) = 26\text{ cm}$.
	}
\end{bt}

% nguồn: De_AI_Toan8_cuoi_hoc_ki_i_1 (nhập 21/11/2026)
\begin{bt}%[8H3H12-1]%[ma:c8e9e5ee]
	Cho hình bình hành $ABCD$ có chu vi bằng $40\text{ cm}$ và $AB - BC = 4\text{ cm}$. Tính độ dài bốn cạnh của hình bình hành $ABCD$.
	\loigiai{
		Nửa chu vi hình bình hành $ABCD$ là $AB + BC = \dfrac{40}{2} = 20\text{ cm}$. \\
Kết hợp với $AB - BC = 4\text{ cm}$, ta tính được: \\
$AB = \dfrac{20 + 4}{2} = 12\text{ cm}$ và $BC = 20 - 12 = 8\text{ cm}$. \\
Vì các cạnh đối của hình bình hành bằng nhau nên $CD = AB = 12\text{ cm}$ và $AD = BC = 8\text{ cm}$.
	}
\end{bt}

% nguồn: De_AI_Toan8_cuoi_hoc_ki_i_1 (nhập 21/11/2026)
\begin{bt}%[8H3H12-1]%[ma:0029b234]
	Cho hình bình hành $ABCD$ có $AB = 2AD$. Gọi $M$, $N$ lần lượt là trung điểm của các cạnh $AB$ và $CD$.
a) Chứng minh tứ giác $AMND$ là hình bình hành.
b) Chứng minh $DM$ là tia phân giác của góc $ADC$.
	\loigiai{
		a) Vì $ABCD$ là hình bình hành nên $AB \parallel CD$ và $AB = CD$. \\
Vì $M$, $N$ là trung điểm của $AB$, $CD$ nên $AM = \dfrac{1}{2}AB = \dfrac{1}{2}CD = DN$ và $AM \parallel DN$. \\
Do đó tứ giác $AMND$ là hình bình hành. \\
b) Ta có $AM = \dfrac{1}{2}AB = AD$, nên tam giác $ADM$ cân tại $A$, suy ra $\widehat{ADM} = \widehat{AMD}$. \\
Mặt khác $AB \parallel CD$ nên $\widehat{AMD} = \widehat{MDC}$ (so le trong). \\
Suy ra $\widehat{ADM} = \widehat{MDC}$, tức $DM$ là tia phân giác của góc $ADC$.
	}
\end{bt}

% nguồn: chuyen_dang (nhập 22/11/2026)
\begin{bt}%[8H3V12-1]%[ma:50b3d0d4]
	
	Chứng minh rằng tổng bình phương hai đường chéo của hình bình hành bằng tổng bình phương bốn cạnh của nó.
	\loigiai{
		\begin{center}
			\begin{tikzpicture}[scale=0.9,font=\footnotesize,line join=round,line cap=round,>=stealth]	
				\def\d{4}
				\def\r{3}
				\path (0:0) coordinate (B)
				++(0:\d) coordinate (C)
				++(50:\r) coordinate (D)
				($(B)+(D)-(C)$) coordinate (A)
				($(A)!0.5!(C)$) coordinate (O);
				\draw (A)--(B)--(C)--(D)--cycle (A)--(C) (B)--(D);
				\foreach \x/ \goc in {A/180,B/180,C/0,D/0,O/90} 
				\fill (\x) circle (1pt)	
				($(\x)+(\goc:3mm)$) node {$\x$};
			\end{tikzpicture}
		\end{center}
		Xét hình bình hành $ABCD$. Gọi $O$ là giao điểm của $AC$ và $BD$ thì $AO$ là trung tuyến của tam giác $ABD$.\\
		Ta có\\
		$AO^2=\dfrac{AB^2+AD^2}{2}-\dfrac{BD^2}{4}\Leftrightarrow \dfrac{AC^2}{4}=\dfrac{AB^2+AD^2}{2}-\dfrac{BD^2}{4}$ $\Leftrightarrow AC^2+BD^2=2\left( AB^2+AD^2 \right)$.}
\end{bt}

% nguồn: chuyen_dang (nhập 22/11/2026)
\begin{ex}%[8H3H12-1]%[ma:c2313e59]
	Trong mặt phẳng $Oxy$, cho hình bình hành $ABCD$ biết $A(2; 1)$, $B(2;-1)$, $C(-2;-3)$. Tọa độ giao điểm hai đường chéo của hình bình hành  $ABCD$ là
	\choice 
		 {$(2; 0)$}
		 {$(2; 2)$}
		{$(0;-2)$} 
		{\True $(0;-1)$} 
	\loigiai{ 
		Giao điểm hai đường chéo $ I $ là trung điểm  của $A C$.
		$$ 
		\heva{&x_I = \dfrac{x_A+x_B}{2}=0 \\ &y_I =  \dfrac{y_A+y_B}{2}=-1.}
		$$ 
		Vậy tọa độ giao điểm hai đường  chéo là $I(0;-1)$.
	} 
\end{ex}

% nguồn: chuyen_dang (nhập 22/11/2026)
\begin{bt}%[8H3V12-1]%[ma:7a8b9c0d]
% nguon: AI tự sinh (Ôn thi HK)
Trong mặt phẳng $Oxy$, cho hình bình hành $ABCD$ có $A(1;2)$, $B(4;1)$, tâm $I\left(\dfrac{9}{2};\dfrac{5}{2}\right)$. Tính góc $\widehat{BAC}$ (làm tròn đến độ).
\shortans{27}
	\loigiai{
$I$ là trung điểm $AC$ nên $C=2I-A=(9-1;5-2)=(8;3)$.\\
$\overrightarrow{AB}=(3;-1)$, $\overrightarrow{AC}=(7;1)$.\\
$\cos\widehat{BAC}=\dfrac{\overrightarrow{AB}\cdot\overrightarrow{AC}}{AB\cdot AC}=\dfrac{21-1}{\sqrt{10}\cdot\sqrt{50}}=\dfrac{20}{\sqrt{500}}=\dfrac{20}{10\sqrt5}=\dfrac{2}{\sqrt5}$.\\
$\widehat{BAC}=\arccos\dfrac{2}{\sqrt5}\approx26{,}57^\circ\approx27^\circ$.
}

\end{bt}

% nguồn: chuyen_dang (nhập 22/11/2026)
\begin{bt}%[8H3H12-1]%[ma:b0d34620]
	Trong mặt phẳng  toạ độ $Oxy$, cho hình bình hành $ABCD$ có diện tích bằng $2$. Biết $A(0;2),B(3;0)$ và giao điểm $I$ của hai đường chéo hình bình hành nằm trên đường thẳng $y=-x$. Tìm hoành độ đỉnh $D$, biết $x_D>-14$.
	\loigiai{ 
		Phương trình đường thẳng $AB$ là $\dfrac{x}{3}+\dfrac{y}{2}=1\Leftrightarrow  2x+3y-6=0$.\\
		Điểm $I$ nằm trên đường thẳng $y=-x$ nên giả  sử $I(t;-t)$.\\
		 Vì $I$ là trung điểm của $AC$ nên $C(2t;-2t-2),I$ là trung điểm của $BD$ nên $D(2t-3;-2t)$.\\
		Ta có  $AB=\sqrt{(3-0)^2+(0-2)^2}=\sqrt{13}$.\\
		Suy ra  $\mathrm{d}(C,AB)=\dfrac{2}{\sqrt{13}}$.\\
		Khi đó $\dfrac{|2\cdot 2t+3(-2t-2)-6|}{\sqrt{2^2+3^2}}=\dfrac{2}{\sqrt{13}} \Leftrightarrow|-2t-12|=2$.\\ 
		Suy ra  $t=-5$ hoặc $t=-7$.
		\begin{itemize}
			\item Với $t=-5$, ta có  $C(-10;8),D(-13;10)$ nhận vì $x_D=-13>-14$.
			\item Với $t=-7$, ta có $C(-14;12),D(-17;14)$ loại  vì $x_D=-17<-14$.
		\end{itemize}
	} 
\end{bt}

% nguồn: chuyen_dang (nhập 22/11/2026)
\begin{bt}%[8H3H12-1]%[Tex hóa SGK CD-CT,T7/22, Lê Minh Thiện Anh]%[ma:af6a5b57]
	Cho ba điểm $A(1;1)$, $B(2;4)$, $C(4;4)$.
	\begin{enumerate}
		\item Tìm tọa độ điểm $D$ sao cho $ABCD$ là một hình bình 	hành.
		\item Tìm tọa độ giao điểm hai đường chéo của hình bình hành $ABCD$.
	\end{enumerate}
	\loigiai{
		\begin{enumerate}
			\item $ABCD$ là hình bình hành $\Leftrightarrow \overrightarrow{AB}=\overrightarrow{DC}$.\\
			Đặt toạ độ của $D(x_D;y_D)$, ta có $\overrightarrow{AB}=\overrightarrow{DC} \Leftrightarrow \heva{&2-1=4-x_D\\&4-1=4-y_D} \Leftrightarrow \heva{&x_D=3\\&y_D=1.}$\\
			Vậy $D(3;1)$.
			\item Gọi $M$ là giao điểm hai đường chéo của hình bình hành $ABCD$.\\
			Ta có $M$ là trung điểm $AC$, suy ra $M(2{,}5;2{,}5)$.
		\end{enumerate}
	}
\end{bt}

% nguồn: chuyen_dang (nhập 22/11/2026)
\begin{bt}%[8H3C12-1]%[ma:ff810c2b]
	Trong mặt phẳng tọa độ $Oxy$, cho hình bình  hành $ABCD$ có $A(4;3)$, $C(8;1)$. Gọi $M$ là trung điểm của cạnh $BC$, $N$ là giao điểm của $BD$ và $AM$. Tính tổng hoành độ của các đỉnh còn lại của hình bình hành $ABCD$, biết $N\left(\dfrac{13}{3};2\right)$.
	\loigiai{ 
		\begin{center}
			\begin{tikzpicture}[scale=1,font=\footnotesize,line join=round,line cap=round,>=stealth]
					 \def\a{6}
					\coordinate[label=left:\scriptsize $B$]  (B) at (0,0);
					\coordinate[label=right:\scriptsize $C$]  (C) at (\a,0);
					\coordinate[label=above:\scriptsize $A$] (A) at ($(B)+(55:.45*\a)$); 
					\coordinate[label=above:\scriptsize $D$]  (D) at ($(C)-(B)+(A)$);
					\coordinate[label=above:\scriptsize $I$] (I) at  ($(C)!.5!(A)$);
					\coordinate[label=below:\scriptsize  $M$] (M) at ($(C)!.5!(B)$);
					\coordinate[label=above:\scriptsize $N$] (N) at  ($(M)!1/3!(A)$);
					 \draw (M)--(A)--(B)--(C)--(D)--(A)--(C)
						 (D)--(B);
					\foreach  \diem in {A,B,C,D,M,N,I}\fill (\diem) circle (1pt);
				\end{tikzpicture}
		\end{center}
		Vì $I$ là tâm  của hình bình hành $ABCD$.\\
		Nên $I$ là trung điểm của  $AC$ suy ra $\heva{&x_I=\dfrac{3+8}{2}=\dfrac{11}{2}\\&y_I=\dfrac{4+1}{2}=\dfrac{5}{2}}\Rightarrow I\left(\dfrac{11}{2};\dfrac{5}{2}\right)$.\\
		Xét tam giác $ABC$ thì $BI$, $AM$ là hai đường trung tuyến nên $N$ là trọng tâm tam giác $ABC$.\\ 
		 Do đó $\heva{&\dfrac{13}{3}=\dfrac{3+x_B+8}{3}\\&2=\dfrac{4+y_B+1}{3}}\Leftrightarrow\heva{&x_B=2\\&y_B=1}\Rightarrow B(2;1)$.\\
		Gọi $D(x_D;y_D)$. Do  $I$ trung điểm của $BD\Rightarrow\heva{&2+x_D=11\\&1+y_D=5}\Leftrightarrow\heva{&x_D=9\\&y_D=4}$
		nên  $D(9;4)$.\\
		Vậy tổng hoành độ của các đỉnh còn lại của hình bình hành  $ABCD$ bằng $11$.
	} 
\end{bt}
